\subsection*{一步一步建立直观}

\paragraph{1. 从导数的误差返回函数的误差}

若导数只用常数近似，函数只得到一条切线；若把导数的线性变化也纳入，积累起来就得到二次修正。积分让一个关于斜率的局部误差变成面积，所以泰勒系数可以通过反复累计来理解，而不必只是代数配系数。



\paragraph{2. 反复积分为什么带来阶乘}

在从零到 $x$ 的区间上，常数积成 $x$，线性项积成 $x^2/2$，再次累计又引入新的分母。右图的指数函数各阶导数重复自身，这让 $1/n!$ 的结构特别明显。每一层累计都把先前的变化分摊到整段区间上。



\paragraph{3. 余项就是尚未解释的累积}

把观察点移远，查看新增高阶项仍未覆盖的差距。积分形式的余项把它写成高阶导数与权重的积分，因而可以用导数上界控制。图形里的两条线看似重合，只表示当前显示尺度下误差小；余项估计才说明这种小有多可靠。


